\subsection{Enchentes históricas}
\lipsum[1-10]

% Décadas de 1940 a 1960
% 1. Enchente de 1941: Uma das mais devastadoras, especialmente em Porto Alegre. 2. Enchente de 1961: Atingiu
%    severamente Porto Alegre e outras regiões. 3. Enchente de 1967: Afetou Uruguaiana e outras cidades ao longo do rio
%    Uruguai.

% Décadas de 1970 a 1990 4 .Enchente de 1973: Inundações significativas em várias regiões do estado. 5. Enchente de
% 1983: Afetou diversas regiões, especialmente áreas ao longo do rio Jacuí. 6. Enchente de 1984: Novamente, grandes
% inundações no rio Jacuí. 7. Enchente de 1995: Inundações em São Leopoldo e no Vale dos Sinos.

% Décadas de 2000 a 2010
% 8. Enchente de 2002: Uma das piores do século, afetando mais de 200 municípios. 9. Enchente de 2007: Grandes chuvas
%    resultaram em inundações severas. 10. Enchente de 2011: Fortes chuvas causaram destruição em cidades como São
%    Lourenço do Sul.

% Décadas de 2010 a 2020 11.1. Enchente de 2011: Em janeiro de 2011, chuvas intensas causaram enchentes que afetaram
% cidades como São Lourenço do Sul, com dezenas de mortos e milhares de desabrigados. 12. Enchente de 2015: Inundações
% generalizadas em julho, afetando Porto Alegre e outras cidades. 13.Enchente de 2017: Chuvas intensas em junho causaram
% inundações em São Sebastião do Caí e Montenegro. Enchente de 2019: Afetou diversas regiões, incluindo a capital, Porto
% Alegre.

% Anos 2020 em diante
% 14. Enchente de 2020: Fortes chuvas causaram inundações em várias regiões. 15. Enchente de 2023: Em setembro, chuvas
%     intensas afetaram mais de 106 municípios e 359 mil pessoas. 16. Enchente de 2024: As enchentes de 2024 foram ainda
%     mais severas, com precipitações recordes entre o final de abril e início de maio, resultando em chuvas de até 700
%     mm em alguns locais. Estes eventos climáticos causaram inundações em 431 dos 497 municípios do estado. Em Porto
%     Alegre, o nível do Lago Guaíba atingiu 5,31 metros, superando o recorde anterior e inundando grandes áreas da
%     cidade. Essas enchentes foram consideradas as piores do Brasil em mais de 80 anos, deslocando milhares de pessoas
%     e causando prejuízos significativos à infraestrutura e economia local . 

\subsection{Enchentes de 2024}
\lipsum[1-6]

% Vale mensionar as enchentes de 2023 como introdução -> Talvez previsões de aumentos de enchentes na região sudeste do
% país? Tem um estudo que foi realizado no governo da Dilma;

% https://www.ncei.noaa.gov/monitoring-content/sotc/global/2024/apr/Brazil-Floods-RioGrandeSul-202404.pdf
% https://web.archive.org/web/20240930134441/https://www.nytimes.com/2024/05/02/world/americas/brazil-rain-floods.html
% https://www.bbc.com/news/world-latin-america-68935249
% https://www.reuters.com/graphics/BRAZIL-RAINS/FLOODS/egpbazlzbvq/

% Cronologia: https://www.bbc.com/portuguese/articles/cd1qwpg3z77o

% As enchentes em Porto Alegre exibiram características tanto de enchentes pluviais quanto fluviais, mas foi o último
% que gerou a maior parte da crise -> Foco no Guaíba. Isso se dever por ter tido a característica de uma enchente
% fluvial, isto é, o Rio foi aquilo que transbordou, mas também pode ser caracterizar como uma enchente pluvial urbana,
% no sentido de que o escoamento do Rio que levou ao aumento do Nível da Água e a falha do escoamento urbano que levou a
% água para dentro da cidade. 

% 5. Climate change, El Niño and infrastructure failures behind massive floods in southern Brazil